\begin{frame}{Selektives Blurring}
\begin{itemize}
    \item Neben ``alle Gesichter blurren'' gibt es einen Lauf fuer konkrete Zielgesichter.
    \item Ein Preview-Lauf zeigt Boxen, Raenge und Track-IDs.
    \item Danach koennen \texttt{--target-track-ids}, \texttt{--target-ranks} oder \texttt{--target-regions} gesetzt werden.
    \item So lassen sich einzelne Personen gezielt anonymisieren, ohne jedes Gesicht im Bild zu verdecken.
\end{itemize}
\vspace{0.5em}
\begin{mynote}[Kontrolle]
Selektives Blurring ist praktisch, braucht aber eine visuelle Pruefung der Track-Stabilitaet.
\end{mynote}
\end{frame}
